\section{Selecting Document Category}
\label{sec-get-labels}

Given a document that the user has bookmarked, the bookmark
tree needs to place the document in proper categories.
To find out to which categories a document belongs,
the traditional approach trains a document classifier
beforehand and applies the classification rules to
classify the document. Taking advantage of that some
search engines can return the document categories, we query
the search engines for the categories,
rather than training our own document classifier.

After a keyword search is issued, a search engine
returns a number of documents. These documents may
belong to different subjects. The user selects some
documents to browse.  Eventually, she/he may save 
a document as a bookmark.

For each bookmark, some search engines return the
labels of the categories to which the bookmark belongs. We can use these
returned labels to classify the bookmark in the bookmark tree.
{\it YAHOO!}~ search engine, for instance,
returns the following three category labels with 
a document titled {\it ``sports car service''}:
\begin{itemize}
\item ``Business and Economy:Company:Automative:Parts:Makes and Models:Saab''
\item ``Business and Economy:Company:Automative:Dealers:Used'', and
\item ``Regional:U.S. States:Delaware:Cites:Wilmington:Business and Shopping".
\end{itemize}

To choose the proper labels to represent the bookmark, we follow four
steps:
\begin{enumerate} 
\item We retain only up to $h - 2$ (excluding root and leaf
level of the tree) prefix of the label.
Since we only have $h - 2$ levels of internal vertices, the 
first $h - 2$ label prefix is sufficient to identify any
directory (and hence category). For instance, suppose $h = 5$.
The labels above become
 \begin{itemize}
 \item ``Business and Economy:Company:Automative'',
 \item ``Business and Economy:Company:Automative'', and
 \item ``Regional:U.S. States:Delaware''.
 \end{itemize}
\item We merge the labels with the same name and assign
each resulting label a weight depending on the fraction of raw
labels it represents. Following the same example, we assign
label ``Business and Economy:Company:Automative'' with 
a weight $\frac{2}{3}$ and 
``Regional:U.S. States:Delaware'' a weight $\frac{1}{3}$. 
\item We check if the weight of a label exceeds the minimum
coverage requirement $c_{min}$. For instance, if we 
enforce that a label must at least cover $10\%$ of the
raw labels, then in this example both labels are
acceptable. Otherwise we may further merge the
labels until all the resulting labels exceed the coverage
threshold. For the lables which coverage are below
$c_{min}$ and cannot be merged with other labels,
we eliminate the labels since they are insignificant.
\item We select the final set of labels to represent the
bookmark. The final set of labels must have a total 
coverage of $g_{min}$. Suppose $g_{min} = 60\%$.
In this example we retain only the 
``Business and Economy:Company:Automative'' 
label since it alone already exceeds the coverage threshold.
If $g_{min} > \frac{2}{3}$, we have to keep both labels for
the bookmark.
\end{enumerate}

\section{Organizing Bookmarks}
\label{sec-build-tree}

The bookmark tree is maintained differently from the
traditional tree structures (e.g., B-trees, R-trees, etc.)
in at least three respects.
First, the shape of the bookmark tree should be stable most of the time.
Structural stability is a strong usability requirement to the bookmark
tree since the users may be confused if the tree is reorganized very often.
Second, the shape of the tree should be adaptive to individual
user's usability preference. A user who prefers generality 
may want to have a less bushy tree with smaller
height; a user who prefers specificity may like
a tree with greater height. (The parameter
$\alpha$ that decides the maximum degree of the vertices
can be tuned to change the shape of the tree. We
discuss the use of $\alpha$ in detail in Sections~\ref{sec-insert}
and~\ref{sec-person}.)
The third difference is that a bookmark may be stored in 
more than one vertex because it has more than one label.
We show how a category
tree handles bookmark insertion, deletion and
update, and how a tree can be personalized to accommodate individual 
preferences in the remainder of this section.

\subsection{Inserting a Bookmark}
\label{sec-insert}

The insertion is performed in two steps. In the
first step, we place the inserted bookmark in the
directories which it belongs to.  In the the second step
we split directories, if the splitting conditions (described shortly)
are met.

When a bookmark is added, we use the bookmark's 
category labels to locate in the tree the deepest vertices 
that match the labels to insert the bookmark.
Suppose Figure~\ref{fig-labels} is our bookmark tree
and we want to insert a bookmark that has three labels ``A:D'',
``C:F'' and ``C:G''. We place the bookmark in the 
subdirectory D under directory A (that matches
label ``A:D'') and in the directory C (that has the 
maximum match for labels ``C:F'' and ``C:G'').

Next, we decide if we should split the directories in which
the bookmark is placed. If the degree of a directory exceeds
the maximum allowable degree $\alpha$, 
the directory is a candidate to be split.
We only split a directory, however, if the split 
decreases the number of children for that directory.
\footnote{A split may not reduce the number of children.
For example, if a new subdirectory can store only one bookmark, the
number of children remains unchanged for the parent directory.}
Thus, we can write the split conditions a
\begin{itemize}
\item $V.d > \alpha$, and
\item $V.d$ must be reduced after the split.
\end{itemize}

When the split conditions are met, we take a greedy approach to 
minimize the number of children in the directory.
Our goal is to make a large move occasionally
rather than making many small moves frequently. 
For all bookmarks in the directory 
that can be grouped and moved, we create subdirectories 
and move them. To prevent frequent split, the splitting
threshold $\alpha$ should be chosen carefully.
We suggest to set the $\alpha$ value slightly larger
than the maximum number of possible subcategories under a category.
For instance, {\it YAHOO}~ classifies documents into 14 top-level
categories. Setting $\alpha$ to, say, $16$ ensures that
no directory split occurs until after the insertion of the $17^{th}$ bookmark.
The first split will create a number of subdirectories
to stage some existing bookmarks (i.e., reducing the number
of bookmarks in the root directory). Furthermore, 
the new subdirectories also provide rooms for staging 
bookmarks. The need
to split the tree again is not likely to happen in
the immediate future.

Finally, we add an alias record to the alias table for
the bookmark if it has multiple labels.

\subsection{Deleting a Bookmark}

When a bookmark is deleted, we 
look up its aliases in the alias table, remove
all its copies from the tree, and delete the
bookmark's alias record from the alias table.
A deletion reduces the number of children 
in the bookmark's parent directories. To keep
the tree stable, however, we do not
collapse a directory even when the
number of children in it becomes very small. 
A directory is removed only when it becomes empty.

\subsection{Update a Bookmark}

\begin{figure}[t]
\epsfysize=180pt
\epsfxsize=320pt
\centerline{\epsffile{fig-bm2.eps}}
\caption{An Update Example}
\label{fig-update}
\end{figure}

The categories to which a bookmark belongs can be changed.
For instance, the content of a document may be changed and thereby
causes it to be reclassified or the search engine may 
reclassify the document into some different categories.
An update can simply be implemented as
an insertion after a deletion. However, the user
should be notified where the bookmarks are moved to.
Whenever a bookmark is moved, we place a forwarding
icon at the old location. 
After the user access the bookmark via the
new path for a few times, the forwarding icon is
removed. Figure~\ref{fig-update} shows that a bookmark
is moved from directories B:E and C (two copies)
to directory B. Both the original directories keep
a forwarding address (in the ovals) referring
to the new directory B.

\subsection{Personalization}
\label{sec-person}

Many personalization features can be support. Here,
we present three ideas and show how they can be implemented:
\begin{enumerate}
\item Hot links:
A hot link is a bookmark identified by the user
as important or is accessed very frequently.
The bookmark tree can color these hot
bookmarks, promote them
to the root directory, or create a hot-link folder
to stage them.

\item Customized categorization:
The user can override the categories that the bookmark
organizer chooses to store the bookmarks. For instance,
the ``Starr Report'' should be categorized into
``Government:US Government:Politics'' but
the user may want to shelve it under the 
``News Event'' folder. In this case, the bookmark
tree uses the user-defined label.
The bookmark tree does not
move the bookmarks that are manually classified by
the user. 

\item Generality versus specificity:
Some people may prefer to keep all files in one directory
while others may prefer to organize files 
in a deep hierarchical tree. We can  
adapt to the user's preference by
changing the value of $\alpha$. A larger $\alpha$ value
makes the tree shallower and hence the categories are more general.
Conversely, a smaller $\alpha$ value makes the tree deeper
and hence the categories can be more specific. 
\end{enumerate}

\subsection{Example}
\label{sec-ex1}

We need a few screen dumps here!

\subsection{Specifications}

We organize bookmarks into a directory tree according to 
bookmarks' labels.  Each directory may have a number of 
bookmarks and child directories that contain bookmarks.  We 
record four pieces of information in each directory vertex:
\begin{itemize}
\item Its label ($V.L$).
\item Its degree ($V.d$): the number of direct children, including 
bookmarks and child directories it has.
\item Its depth in the tree ($V.\delta$). 
\item Its children set ($V.v$), including both directories and bookmarks.
\end{itemize}
A leaf vertex is a bookmark vertex, so its children set is always empty
($V.v = \emptyset$).  For a leaf vertex, we record the name of the 
bookmark in $V.B$.

Let $L_s$ denote the set of category labels of a bookmark $B$.
Let $L$ denote an individual category label, which consists of
a list of sublabels concatenated by colons. 
Figures~\ref{fig-insert} depicts the
algorithms to insert a bookmark.
For each label of the bookmark, 
algorithm {\it Bookmark-Insert} first calls procedure 
{\it Find-Max-Matching-Vertex} (depicted in Figure~\ref{fig-max})
to find the deepest directory
to insert the bookmark. After updating some bookkeeping
information for the inserted bookmark and directory, 
it calls procedure {\it Split-Vertex}
to split the directory if the directory's degree exceeds the
splitting threshold.  The {\it Split-Vertex} procedure (not
documented) checks if splitting reduces the number of children for
the directory. If so, it creates subdirectories, moves
bookmarks, and update bookkeeping parameters for
bookmarks and directories involved.
Finally, procedure {\it Bookmark-Insert}
updates the alias table for the bookmark,
if necessary. 

Algorithm {\it Bookmark-Delete} in Figure~\ref{fig-delete}
describes the deletion steps. The algorithm first
gets all the aliases of the to-be-deleted bookmark from the alias
table. For each alias, it finds the directory by calling
procedure {\it Find-Max-Matching-Vertex} and then removes
the bookmark from that directory. Procedure 
{\it Collapse-Vertex} (not documented) is invoked
if the degree of the vertex, after removing the bookmark,
becomes zero. The procedure may recursively remove a
number of vertices if they also become empty.

\begin{figure}
{\bf Bookmark-Insert($T,L_s,B$)}
\begin{nlist}
\item for each $L \in L_s$
 \begin{nlist}
 \item $V^* \leftarrow$ {\bf Find-Max-Matching-Vertex}$(T,L)$;
 \item new $V$; /* allocate a new vertex */
 \item $V.B \leftarrow B$; 
 \item $V.L \leftarrow L$;
 \item $V.\delta \leftarrow V^*.\delta + 1$;
 \item $V^*.v \leftarrow V^*.v \cup V$;
 \item $V^*.d \leftarrow V^*.d + 1$;
 \item if ($V^*.d > \alpha$

       {\bf Split-Vertex}$(V^*)$;
 \end{nlist}
 \item if ($|L_s| > 1$) {\bf Update-Alias-Table}$(AT, L_s, B)$;
\end{nlist}
\caption{Algorithm Bookmark Insert}
\label{fig-insert}
\end{figure}

\begin{figure}
{\bf Bookmark-Delete($T,B$)}
\begin{nlist}
\item $L_s \leftarrow$ {\bf Find-Aliases}$(AT, B)$;
\item for each $L \in L_s$
 \begin{nlist}
 \item $V^* \leftarrow$ {\bf Find-Max-Matching-Vertex}$(T,L)$;
 \item for each $V \in V^*.v$
  \begin{nlist}
  \item if ($V.B = B$) 
   \begin{nlist}
   \item $V^*.v \leftarrow V^*.v - V$;
   \item delete $V$;
   \end{nlist}
  \item $V^*.d \leftarrow V^*.d - 1$;
  \item if ($V^*.d = 0$) {\bf Collapse-Vertex}$(T, V^*)$;
  \end{nlist}
 \end{nlist}
\end{nlist}
\caption{Algorithm Bookmark Delete}
\label{fig-delete}
\end{figure}

\begin{figure}
{\bf Find-Max-Matching-Vertex(T,L)}
\begin{nlist}
\item $found \leftarrow false$;
\item $l \leftarrow head[L]$;
\item $V \leftarrow root[T]$;
\item $V^* \leftarrow root[T]$;
\item while ($l \ne nil$) $\{$
 \begin{nlist}
 \item for each $V \in V.v$ $\{$
  \begin{nlist}
  \item if ($l = V.L$) $\{$ /* matching the label */
   \begin{nlist}
   \item $V^* \leftarrow V$; \item $found \leftarrow true$;
   \item break; $\}$
   \end{nlist}
  \item $\}$
  \end{nlist}
 \item if (found) /* found a match, next level */
  \begin{nlist}
  \item $l \leftarrow next[l]$; 
  \item $found \leftarrow false$;
  \end{nlist}
 \item else break; /* ready to return */
 \end{nlist}
\item $\}$ /* while */
\item return($V^*$);
\end{nlist}
\caption{Procedure Find-Max-Matching-Vertex}
\label{fig-max}
\end{figure}


